% Generated draft: requires mathematical and source review.
\documentclass[11pt]{article}
\usepackage[margin=25mm]{geometry}
\usepackage{fontspec}
\setmainfont{Latin Modern Roman}
\usepackage{amsmath,amssymb,mathtools,unicode-math}
\setmathfont{Latin Modern Math}
\usepackage{enumitem,needspace,xurl,hyperref}
\setlength{\parindent}{0pt}
\setlength{\parskip}{5pt}
\setlength{\emergencystretch}{3em}
\newcommand{\N}{\mathbb{N}}
\newcommand{\Z}{\mathbb{Z}}
\newcommand{\Q}{\mathbb{Q}}
\newcommand{\R}{\mathbb{R}}
\newcommand{\C}{\mathbb{C}}
\newcommand{\F}{\mathbb{F}}
\newcommand{\A}{\mathbb{A}}
\newcommand{\PP}{\mathbb{P}}
\newcommand{\E}{\mathbb{E}}
\newcommand{\eps}{\varepsilon}
\newcommand{\dd}{\,\mathrm{d}}
\begin{document}
\section*{UnsolvedMath ID: 800452}
\textbf{Source problem code:} catalogue entry
\subsection{Definitions and mathematical statement}
Let \(n\geq 1\) be an integer and let \(\F_2=\{0,1\}\) be the field with two elements. For \(x,y\in\F_2^n\), define the inner product modulo \(2\) function by
\[
\operatorname{IP}_{2n}(x,y)=(-1)^{\sum_{i=1}^{n}x_i y_i},
\]
where the exponent is interpreted modulo \(2\).

The supplied extract does not explicitly specify the monomial basis used to measure PTF sparsity. The following formulation uses signed coordinates: encode a bit \(u_j\in\F_2\) by \(z_j=(-1)^{u_j}\in\{-1,1\}\). For a Boolean function \(F\colon\{-1,1\}^m\to\{-1,1\}\), a polynomial threshold function (PTF) representation means a real multilinear polynomial
\[
p(z)=\sum_{S\subseteq\{1,\ldots,m\}}a_S\prod_{j\in S}z_j
\quad\text{such that}\quad
F(z)p(z)>0\quad\text{for every }z\in\{-1,1\}^m.
\]
The empty product is \(1\), so a nonzero constant coefficient counts as one monomial. There is no degree restriction. Define
\[
\operatorname{sp}_{\mathrm{PTF}}(F)
=\min_p\bigl|\{S:a_S\ne0\}\bigr|,
\]
where the minimum is over these representations. This convention requires nonzero polynomial values at every input.

Define the signed-coordinate version of the inner product function by
\[
\widetilde{\operatorname{IP}}_{2n}
\bigl((-1)^{x_1},\ldots,(-1)^{x_n},
      (-1)^{y_1},\ldots,(-1)^{y_n}\bigr)
=\operatorname{IP}_{2n}(x,y).
\]
Under the stated convention, the source question is whether
\[
\operatorname{sp}_{\mathrm{PTF}}
\bigl(\widetilde{\operatorname{IP}}_{2n}\bigr)\geq 3^n
\qquad\text{for every integer }n\geq1.
\]
Equivalently, must every such PTF representation have at least \(3^n\) nonzero monomials? Ordinary monomials in the original \(0,1\) coordinates give a different sparsity notion; the intended basis must therefore be confirmed against the source before this conditional formulation is accepted.

The source additionally records a much stronger question independently asked by Rocco Servedio. With the same PTF convention, for every integer \(n\geq1\), every pair of functions \(f,g\colon\{-1,1\}^n\to\{-1,1\}\), and all positive integers \(s,t\), does
\[
\operatorname{sp}_{\mathrm{PTF}}(f)\geq s,
\qquad
\operatorname{sp}_{\mathrm{PTF}}(g)\geq t
\quad\Longrightarrow\quad
\operatorname{sp}_{\mathrm{PTF}}(f\oplus g)\geq st
\]
hold, where \(f\oplus g\colon\{-1,1\}^{2n}\to\{-1,1\}\) is defined on disjoint inputs by \((f\oplus g)(z,w)=f(z)g(w)\)?
\subsection{Short English statement}
The question asks whether every polynomial threshold representation of inner product modulo \(2\) on \(2n\) bits needs at least \(3^n\) monomials. The formulation above measures monomials in signed coordinates; the supplied extract leaves this basis convention unresolved. A stronger question attributed to Rocco Servedio asks whether PTF sparsity lower bounds multiply when two Boolean functions on disjoint inputs are multiplied.
\subsection{Sources}
\begin{itemize}
\item Local UnsolvedMath registry, permanent ID \(800452\), exact source problem code \texttt{catalogue entry}, title: \emph{Analysis of Boolean Functions --- PTF Sparsity for Inner Product Mod 2}. The supplied registry labels the problem open; this is a supplied status, not an independently verified mathematical or literature assessment.
\item O'Donnell, \emph{Open Problems in Analysis of Boolean Functions} (2012), Problem 23, ``PTF Sparsity for Inner Product Mod 2.'' Supplied source URL: \url{https://arxiv.org/abs/1204.6447}. The record reports access on 2026-10-16 and extraction by \texttt{arxiv\_latex}; these are supplied provenance details, not a claim of independent source access.
\item The supplied background attributes the proposal to Srikanth Srinivasan, 2010. The registry's proposer and year fields instead list O'Donnell and 2012, corresponding to the cited source list; both attributions are retained.
\item The supplied remarks attribute the stronger multiplicative PTF sparsity question to Rocco Servedio, who independently asked it. No separate reference for this question is supplied.
\end{itemize}
\end{document}
